\section{Finite bridges for Bayesian reasoning models}
\label{sec:annex-p19-rationality-4-5}

\begin{annexitem}{annex:p19.r45.coherence}{thm:p19-r45-coherence}{The four cells, sharpness and uncertainty inequality}
\begin{proof}
The antecedent/consequent cells have masses
\[
 P(A\cap B)=ab,\quad P(A\cap B^c)=b(1-a),\quad
 P(A^c\cap B)=z,\quad P(A^c\cap B^c)=1-b-z.
\]
Nonnegativity is equivalent to \(0\leq z\leq1-b\), since
\(a\in[0,1]\) and \(b\in(0,1)\). Thus
\(P(B)=ab+z\) lies in exactly the claimed interval. Conversely,
for any number in that interval choose its difference from \(ab\)
as \(z\). These four nonnegative masses sum to one and give the
prescribed antecedent probability and conditional probability. This
proves attainability of every point, including both ends.

Since \(P(B)\geq ab\), its uncertainty is at most \(1-ab\).
Moreover,
\[
 ((1-a)+(1-b))-(1-ab)=(1-a)(1-b)\geq0,
\]
which proves the general uncertainty-sum property. No selected interior
point or assumption of independence is used.

The material conditional omits exactly \(A\cap B^c\), and so
its probability is \(1-b(1-a)=1-b+ab\). Subtracting \(a\)
gives \((1-a)(1-b)\), proving the equality characterization.
If \(b=1\), the other two cells are null and the consequent has
mass \(a\). If \(b=0\), the ratio defining \(P(B\mid A)\)
has denominator zero and is not assigned a value by ordinary
conditioning. A different conditional-event framework would need its
own declaration; no endpoint convention has entered this proof.
\end{proof}
\end{annexitem}

\begin{annexitem}{annex:p19.r45.jeffrey}{thm:p19-r45-jeffrey}{Mixture, cell preservation and nonidentification}
\begin{proof}
The two positive-event conditional laws under \(P_0\) are probability
measures. A convex mixture with weights \(t,1-t\) is therefore a
probability measure, also at either endpoint. Conditional laws on the
two cells give values 1 and 0 respectively to \(A\), so the mixture
assigns it \(t\).

For an event \(D\),
\[
 P_1(D\cap A)=tP_0(D\mid A),\qquad
 P_1(D\cap A^c)=(1-t)P_0(D\mid A^c).
\]
When \(0<t<1\), dividing by the respective cell probabilities
proves preservation. Conversely, any law with that marginal and those
two conditional laws satisfies precisely these two joint-cell equations.
Adding them proves the mixture formula for every \(D\), hence
uniqueness. At \(t=1\) or 0 the mixture reduces directly to the
corresponding old conditional law. The new zero-weight cell is null;
preserving a version there is not an assertion that its conditional
probability is determined by a ratio.

For the nonidentification boundary, fix any \(u\in[0,1]\).
Give the four \(A,B\) cells masses
\[
 tu,\quad t(1-u),\quad (1-t)u,\quad (1-t)(1-u),
\]
in the order \(AB,AB^c,A^cB,A^cB^c\). They are nonnegative,
sum to one and have marginals \(P_1(A)=t\), \(P_1(B)=u\).
This works at \(t=0,1\) as well as in the interior. Without other
constraints or rigidity, every consequent marginal is attainable.
\end{proof}
\end{annexitem}

\begin{annexitem}{annex:p19.r45.quantifiers}{thm:p19-r45-quantifiers}{Positive mass, existence and a rare-category counterexample}
\begin{proof}
The event \(C\cap B\cap A^c\) is a subset of the stipulated
null event \(B\cap A^c\), so its probability is zero. Decomposing
\(C\cap B\) by \(A,A^c\) yields
\[
 P(C\cap B)=P(C\cap B\cap A)\leq P(C\cap A).
\]
The left side is positive by hypothesis. In particular both subject
categories needed to form the conditional readings are nonnull.

On a finite domain with positive probability at every object, every
nonempty event has positive mass. Consequently a null exception set
is empty and a positive intersection is nonempty. If zero-mass objects
are permitted, use a two-object domain with masses 1 and 0, let \(B\)
contain both and \(A\) only the positive-mass object. Then
\(P(B\cap A^c)=0\), although \(B\) is not a subset of \(A\).
This establishes the need for the stated support condition in the
universal translation, without altering ordinary probability laws.

For the most counterexample, use a uniform domain of one million
objects. Let \(C\) have 100 objects and \(B\) 9,000, with
99 objects in their intersection. Let \(A=B\setminus C\), so
\(A\) has 8,901 objects. The remaining objects have no other
membership constraints. Then
\[
 P(B\mid C)=99/100=.99,\qquad
 P(A\mid B)=8901/9000=.989,\qquad P(A\mid C)=0.
\]
Both first probabilities lie strictly between .95 and 1, but the
conclusion's conditional probability is zero. All three predicates
are rare in this domain as well; a bare rarity assumption would not
remove this particular counterexample.
\end{proof}
\end{annexitem}

\begin{annexitem}{annex:p19.r45.sample}{thm:p19-r45-sample}{What absence from a sample proves}
\begin{proof}
For each independent draw, the noncounterexample event has probability
\(1-r\). Independence makes the probability of their conjunction
the product of these \(n\) probabilities, \((1-r)^n\).
At \(r=.01,n=7\), this is \(.99^7=.93206534790699\).
If \(0<r<1\) and \(n\) is finite, this absence event has
positive probability even though counterexamples exist with positive
mass. It therefore cannot logically establish a universal statement.
The cases \(r=0,1\) yield probabilities 1 and 0 respectively for
\(n\geq1\). The calculation does not estimate a human response
law, nor prove that the model's draws correspond to observed memory
instances.
\end{proof}
\end{annexitem}

\begin{annexitem}{annex:p19.r45.description}{thm:p19-r45-description}{Normalization and the MAP tradeoff}
\begin{proof}
Finiteness, nonemptiness and finite lengths ensure that \(Z\) is finite
and strictly positive. The normalized weights are positive and sum to
one. Let \(m=\sum_h\pi(h)\ell(h)>0\) be the stipulated prior
predictive evidence probability. Bayes gives
\[
 \pi(h\mid D)=\frac{2^{-L(h)}\ell(h)}{Zm}.
\]
The denominator is common to every candidate. For candidates with
positive likelihood, maximizing the numerator is equivalent to
minimizing its negative base-two logarithm, namely
\(L(h)-\log_2\ell(h)\). A zero-likelihood candidate has posterior
zero and infinite objective; at least one candidate has positive
likelihood because \(m>0\). On the finite set a minimizer exists,
and equal values remain ties. For a common positive likelihood, its
logarithm is a constant and minimization reduces to code length.
Otherwise no such reduction follows without a separate ordering
argument.

In the two-candidate example, unnormalized prior weights are
\(1/2\) and \(1/8\), so their normalized priors are \(4/5\)
and \(1/5\). Multiplication by likelihoods \(1/100\) and 1 gives
\(1/125\) and \(1/5\), with total \(26/125\). The longer
candidate's posterior is \((1/5)/(26/125)=25/26\), demonstrating
the likelihood qualification directly. No claim about a uniquely
privileged code language is used.
\end{proof}
\end{annexitem}

\begin{annexitem}{annex:p19.r45.information}{thm:p19-r45-information}{Entropy identity, nonnegativity and realized uncertainty}
\begin{proof}
Write \(J(h,y)=\pi(h)k_a(y\mid h)\) for the joint law and
\(R(h,y)=\pi(h)m_a(y)\) for its marginal-product law. Whenever
\(J(h,y)>0\), the stipulated positive prior and mixture definition
give \(R(h,y)>0\). Expanding conditional posterior entropy and
using \(\pi_a(h\mid y)=J(h,y)/m_a(y)\) gives
\[
 S(\pi)-\sum_y m_a(y)S(\pi_a(\cdot\mid y))
 =\sum_{J>0}J(h,y)\log_2\frac{J(h,y)}{R(h,y)}.
\]
Only outcomes with positive marginal enter the entropy sum. Dividing
\(J/R\) cancels the prior factor and gives the displayed likelihood
ratio \(k_a/m_a\).

For \(x>0\), \(\log x\leq x-1\), with equality exactly at
\(x=1\): differentiating \(x-1-\log x\) shows a global minimum
zero at 1. Applying this to \(R/J\), summing over \(J>0\),
and converting the logarithm base yields
\[
 \sum_{J>0}J\log_2(J/R)
 \geq\frac{1-\sum_{J>0}R}{\log 2}\geq0.
\]
If the sum is zero, both inequalities are equalities. Thus \(R\)
has no mass outside the support of \(J\), and every positive
\(J\) equals \(R\). Conversely, \(J=R\) makes the sum zero.
Since every prior is positive, this equality is precisely
\(k_a(y\mid h)=m_a(y)\) for all hypotheses and outcomes. A
zero marginal implies every associated likelihood is zero as well.

The binary illustration has outcome marginal one-half for both
success and failure; Bayes gives posteriors .8/.2 and .2/.8.
Both posterior entropies equal \(S(.8,.2)\), establishing the stated
gain. To demonstrate the realized-outcome qualification, instead take
prior (.9,.1) and a success likelihood of \(1/9\) under the first
hypothesis and 1 under the second. Success has probability .2 and
posterior (.5,.5), whose entropy 1 exceeds the prior entropy. Failure
has probability .8 and posterior (1,0). Expected posterior entropy is
.2, below the prior entropy, so expected gain is nevertheless positive.

The formula contains neither utility nor acquisition cost. Adding a
large positive test cost can make any finite information-gain test
undesirable under a decision criterion, without changing this identity.
Likewise, the posterior belief difference only has a decision value
after losses/actions are specified. The expectation result alone
therefore establishes neither of the stronger optimization claims.
\end{proof}
\end{annexitem}

\begin{annexitem}{annex:p19.r45.slope}{thm:p19-r45-slope}{Action-specific loss and the baseline}
\begin{proof}
The expected utility of action 1 is \(b-cq_1\), and of action 0
is \(-cq_0\). Subtracting gives \(b-c(q_1-q_0)\).
Its positive, zero and negative values give strict preference, a tie,
and reverse strict preference, respectively. If \(q_0=0\), the
strict comparison reduces to \(b>cq_1\); if it is not zero, that
reduction does not generally hold. The example has difference
\(3-10(.4-.2)=1\), but \(3-10(.4)=-1\).
The stipulated utilities/laws provide the calculation. An observational
association need not equal those prospective action-specific
probabilities, so the calculation establishes no observational
identification or policy forecast.
\end{proof}
\end{annexitem}
